\section{The domestic clause, vote by vote}
\label{sec:F.6}
\textbf{The record, as of 24 September 2026.} The June 2025 strikes on Iran are the pre-2026 instance: no authorization existed, and the Senate's war-powers resolution failed 47 to 53 \autocite{senate2025kaine}, which changes nothing under the clause, since what it reads is the absence, not the vote. The 2026 campaign supplies every marker the clause reads, and none was met. The strikes began on 28 February; the administration reported to Congress under the War Powers Resolution on 2 March, which started the sixty days. The Charter clause's verdict on the same campaign is in \S\ref{sec:6.4}. No declaration was made and no authorization passed: one was introduced in the House on 7 May and never left committee \autocite{hjres176}. On 1 May, the sixtieth day, the President wrote to Congress that the hostilities begun on 28 February had terminated, resting on a ceasefire declared on 7 April, while a naval blockade continued and the Secretary of Defense testified that a ceasefire stops the clock \autocite{cbs2026terminated,lawfare2026sixtydays}. Hostilities resumed in July under no new authorization. On 3 June the House, 215 to 208, and on 23 June the Senate, 50 to 48, adopted a concurrent resolution directing the President to remove forces from hostilities against Iran \autocite{hconres86,npr2026housejune,asil2026iranwpr}; a White House official said it had no significance \autocite{abc2026nosignificance}, and nothing changed. No binding joint resolution came to a vote on passage: the Senate discharged one, 50 to 47, on 19 May \autocite{cbs2026sjres185discharge,lawfare2026resolutions}, refused to take it up, 47 to 50, on 24 June \autocite{senate2026rollcall192}, and a later motion to discharge failed 47 to 49 in July \autocite{thehill2026discharge}. A third concurrent resolution passed the House 220 to 204 on 15 September, and the second, passed by the House in July 214 to 208, failed in the Senate 49 to 50 on 24 September \autocite{npr2026senatesept}. The clause doesn't need the votes. What it reads is the absence of an authorization, and by 1 May that absence was on the record in the President's own letter, which conceded that hostilities had begun and claimed only that they had ended.
